% Pursuing Stacks (cotes 134-2 to 134-8): archive cote 134-2, site page 81 (PDF page 82).
% notes to ch. I. Draft compared: ps1.tex:1385-1387 ; ps1.tex:1394-1406 ; ps1.tex:1597-1600 ; ps1.tex:1621-1621 ; ps1.tex:1623-1648.
% AI-transcribed by Claude Opus 5.5 (claude-opus-5-5, adaptive thinking, effort high), Message Batches API, batch ps1 (msgbatch_01Rv5hCaNTtUt4JHx1yqoHPL), 2026-11-12; not yet checked by eye.
\scanpage{134-2/81}
\mspage{65}

\inserted{a} considerably wider sense than I expected,  when writing this letter to
Larry Breen.  It now appears that the whole theory of stacks of groupoids
will depend on the choice of a ``coherator'' $\underline{C}$ , as seen in section 13.

($^{26}$)\quad Tim Porter pointed out to me that ``Dold Puppe'' is an \inserted{in}accurate name
for this basic theorem,  which should be called \uline{Dold-Kan theorem}.

($^{27}$)\quad As was pointed out to me by Ronnie Brown,  this structure was
already well-known to J.H.C. Whitehead,  under the name of ``crossed modu-%
le\xout{\unsure{s}}'',  and extensive use and extensive generalizations of this notion
(in quite different directions from those I was having in mind,  in
terms of Gr-stacks over an arbitrary topos) have been made by him and
others.  With respect to the question on next page,  of generalizing this
notion of ``non commuative chain complex'' from length one to length two,
Ronnie says there is a work in preparation by D. Conduché `` Modules croi-%
sés généralisés de longueur 2 '' .

($^{28}$)\quad Tim Porter pointed out to me that work on étale tubular neighbour-%
hoods was done by D.A. Cox ; ``Algebraic tubular neighbourhoods I,II',
Math. Scand. 42 (1978)  211-228, 229-242.  I've not seen yet this work,
and can't say therefore whether it meets the rather precise expectations
I have for a theory of tubular neighbourhoods,  for the needs of a dévis-%
sage theory of stratified schemes (or, more generally, stratified topoi).

($^{29}$)\quad This is the typical game \xout{\unsure{to be}} embodied in the ``derivator'' associa\xout{\unsure{te}}%
\hand{-}ted to the theory (Hot) of usual homotopy types (compare section 69).

($^{30}$)\quad Some more details on this program are outlined in \hand{``}\uline{Esquisse d'un
Programme}\hand{''}(section 5),  in \uline{Réflexions} \uline{Mathématiques} 1 .

($^{31}$)\quad I was informed by knowledgeable people soon later that the answer is
well known to be negative,  by working with ``rational homotopy types''
(the cohomology of which is made up with vector spaces over $\underline{Q}$). It is well
known indeed that a 1-connected rational homotopy type is \uline{not} known from
it\unsure{\deleted{'}}s rational cohomology ring alone,  which \deleted{is supposed (I confess I did'nt
check) to} contain\hand{s} already all the information I was contemplating\deleted{\unsure{.(*)}}\inserted{\deleted{\illegible{1}} \deleted{(*)} At least \inserted{this is so} if we assume that \deleted{\unsure{the}} $H^i(X)$ is of finite dimension over $\mathbb{Q}$ for all $i$. But is there a counterexample still when $X$ is a homotopy type ``of finite type'' ?}

($^{32}$)\quad This \inserted[typed]{``problem'' is met with by the notion of a ``derivator''}\xout{is the non \unsure{commutative} version of the similar problem} , which
``was in the air'' already by the late sixties,  but was never developped \deleted{\illegible{1}}
(instead even derived categories became tabu in the seventies ...). Compare
with note ($^{24}$).

\margin{81}
